\documentclass[../DoAn.tex]{subfiles}
\begin{document}
\section{Kết luận}
Đồ án đã tập trung giải quyết bài toán đặt vé và quản lý sự kiện trên thiết bị di động theo hướng xây dựng một hệ thống tích hợp nhiều nghiệp vụ trong cùng một nền tảng. So với các sản phẩm định hướng cộng đồng như Meetup, hệ thống của đồ án chưa đặt trọng tâm vào quy mô cộng đồng rộng hoặc hệ sinh thái người dùng lớn, nhưng đã hướng nhiều hơn tới sự gắn kết giữa khâu khám phá sự kiện, đặt vé, lưu vé điện tử, nhận thông báo và kiểm soát tham dự \cite{meetup2026about}. So với các nghiên cứu học thuật về gợi ý sự kiện hoặc gợi ý địa điểm tổ chức, chẳng hạn các công trình về đề xuất sự kiện trong môi trường mạng xã hội và tối ưu lựa chọn cho nhóm người dùng \cite{cena2014event,liao2020acger,zhang2019gevr}, sản phẩm của đồ án không đi sâu vào mô hình đề xuất thông minh, nhưng lại có ưu điểm ở việc hiện thực hóa một luồng nghiệp vụ tương đối hoàn chỉnh dưới dạng ứng dụng có thể sử dụng thử nghiệm trên thiết bị di động. Nói cách khác, các công trình nghiên cứu tương tự mạnh về phân tích và tối ưu hóa quyết định, trong khi hệ thống của đồ án mạnh hơn ở khía cạnh tích hợp chức năng và hiện thực hóa sản phẩm.

Trong suốt quá trình thực hiện đồ án, kết quả đạt được nổi bật nhất là xây dựng được một hệ thống gồm ứng dụng di động và backend có khả năng vận hành các nghiệp vụ cốt lõi của bài toán. Cụ thể, hệ thống đã hỗ trợ đăng ký và đăng nhập người dùng, quản lý hồ sơ, tìm kiếm và lọc sự kiện, xem chi tiết sự kiện, lưu sự kiện yêu thích, tạo sự kiện, quản lý hạng vé, đặt vé, xác nhận thanh toán, phát hành vé điện tử, sinh mã \gls{QR} cho vé, check-in tại sự kiện, nhắn tin, gửi lời mời, gửi thông báo và nhắc lịch tự động. Ngoài ra, hệ thống còn có cơ chế phân quyền giữa người dùng, người tổ chức và quản trị viên, đồng thời tích hợp chức năng gợi ý nội dung sự kiện bằng AI để hỗ trợ người tổ chức khi tạo sự kiện mới. Những kết quả này cho thấy mục tiêu chính của đề tài là xây dựng một ứng dụng đặt vé sự kiện mobile-first có tính tích hợp đã được đáp ứng ở mức khả thi.

Bên cạnh các kết quả đạt được, đồ án vẫn còn một số hạn chế. Thứ nhất, hệ thống mới dừng ở mức thử nghiệm và chưa tích hợp đầy đủ cổng thanh toán thực tế ở môi trường triển khai thương mại. Thứ hai, các chức năng gợi ý thông minh mới chỉ dừng ở hỗ trợ tạo nội dung sự kiện, chưa phát triển thành cơ chế đề xuất sự kiện cá nhân hóa cho người dùng cuối. Thứ ba, hệ thống chưa đi sâu vào các vấn đề triển khai quy mô lớn như cân bằng tải, giám sát hệ thống, tối ưu hiệu năng cho lượng truy cập rất cao hoặc thống kê phân tích hành vi người dùng trên quy mô thực tế. Ngoài ra, việc kiểm thử hiện nay mới phù hợp ở mức chức năng chính, chưa bao phủ đầy đủ các tình huống kiểm thử tự động và các bài toán vận hành ngoài thực địa.

Qua quá trình thực hiện đồ án, có thể rút ra một số bài học kinh nghiệm quan trọng. Việc phân tách rõ vai trò giữa frontend di động, backend nghiệp vụ và cơ sở dữ liệu giúp cho quá trình phát triển và bảo trì thuận lợi hơn. Bên cạnh đó, việc tổ chức hệ thống theo các mô-đun như xác thực, sự kiện, vé, thanh toán, chat và thông báo cho thấy hiệu quả rõ rệt khi cần bổ sung hoặc điều chỉnh chức năng. Một kinh nghiệm khác là với các bài toán có nhiều vai trò tham gia như người dùng, người tổ chức và quản trị viên, việc mô hình hóa use case và quy trình nghiệp vụ từ sớm có ý nghĩa rất quan trọng trong việc tránh thiếu sót chức năng và giảm xung đột khi triển khai. Cuối cùng, đồ án cho thấy rằng một hệ thống sự kiện hiện đại không chỉ cần giao diện đẹp hay quy trình mua vé đơn giản, mà còn cần tính liên kết chặt chẽ giữa các nghiệp vụ để đem lại trải nghiệm xuyên suốt cho người dùng.

\section{Hướng phát triển}
Trong giai đoạn tiếp theo, trước hết cần tiếp tục hoàn thiện các chức năng đã được xây dựng trong đồ án để hệ thống có thể tiến gần hơn tới mức sẵn sàng triển khai thực tế. Một số công việc cần ưu tiên gồm tích hợp cổng thanh toán thật thay cho luồng xác nhận thanh toán mô phỏng, tăng cường cơ chế xử lý lỗi và xác thực đầu vào ở các API quan trọng, hoàn thiện trải nghiệm giao diện trên nhiều kích thước thiết bị, và nâng cao độ ổn định của các chức năng thông báo đẩy, nhắc lịch và chat thời gian thực. Bên cạnh đó, cần bổ sung kiểm thử tự động cho các nghiệp vụ trọng yếu như tạo sự kiện, đặt vé, thanh toán và check-in để giảm rủi ro khi mở rộng hệ thống. Việc tối ưu hóa phần tìm kiếm, lọc sự kiện, quản lý hình ảnh sự kiện và đồng bộ dữ liệu giữa ứng dụng khách với máy chủ cũng là những hạng mục cần thiết để sản phẩm hoạt động tin cậy hơn trong môi trường thực tế.

Sau khi các chức năng hiện tại được hoàn thiện, hệ thống có thể được mở rộng theo nhiều hướng mới nhằm nâng cao chất lượng và giá trị sử dụng. Hướng thứ nhất là phát triển cơ chế gợi ý sự kiện cá nhân hóa dựa trên lịch sử tương tác, sở thích, vị trí và hành vi tham gia của người dùng, từ đó đưa hệ thống tiến gần hơn tới các hướng nghiên cứu đã được đề cập trong các công trình học thuật. Hướng thứ hai là mở rộng phần quản trị và phân tích dữ liệu, chẳng hạn bổ sung dashboard thống kê cho người tổ chức và quản trị viên để theo dõi lượt quan tâm, số vé đã bán, tỷ lệ check-in, phản hồi của người dùng và hiệu quả của từng sự kiện. Hướng thứ ba là mở rộng khả năng triển khai bằng cách đưa hệ thống lên hạ tầng đám mây, hỗ trợ lưu trữ ảnh tập trung, tăng khả năng mở rộng theo số lượng người dùng và bổ sung cơ chế sao lưu, giám sát, ghi log tập trung. Cuối cùng, hệ thống có thể tiếp tục được phát triển theo hướng đa nền tảng và đa dịch vụ, ví dụ hỗ trợ thêm web quản trị, tích hợp bản đồ và định vị sâu hơn, hỗ trợ đa ngôn ngữ, đăng nhập liên thông hoặc mở rộng sang các mô hình sự kiện có sơ đồ chỗ ngồi, bán vé theo khu vực và quản lý chương trình chi tiết theo thời gian.
\end{document}
